\documentclass[11pt]{article}

\usepackage[margin=1in]{geometry}
\usepackage{hyperref}

\title{DIP-004: Independent Verification of Deterministic Decisions}

\author{Pavan Dev Singh Charak}

\date{2026}

\begin{document}

\maketitle

\begin{abstract}

This document describes the verification model used by the
Decision Integrity Protocol (DIP). The verification process
allows independent systems to validate the integrity of
decision records.

\end{abstract}

\section{Introduction}

Verification is a core feature of the Decision Integrity
Protocol. Independent verification ensures that decision
records remain trustworthy even when produced by automated
systems.

\section{Verification Goals}

The verification process ensures:

\begin{itemize}
\item integrity of decision payloads
\item authenticity of decision signatures
\item consistency of ledger records
\end{itemize}

\section{Verification Architecture}

The verification process involves the following components:

\begin{itemize}
\item decision payload
\item payload signature
\item decision ledger
\item verification engine
\end{itemize}

\section{Verification Procedure}

A verifier validates decision records using the following steps:

\begin{enumerate}

\item Parse the decision payload

\item Reconstruct canonical payload representation

\item Compute the payload hash

\item Verify the cryptographic signature

\item Confirm ledger inclusion

\end{enumerate}

\section{Trust Model}

Verification systems operate independently of the
decision-producing system.

This separation ensures that decision evidence can be
validated without relying on the originating system.

\section{Implementation}

A reference implementation of the verification engine
is available in the DIP repositories.

\section{Conclusion}

Independent verification enables deterministic governance
systems that can be audited and validated by external
parties.

\end{document}